\section{Reduction of full nef dimension to Calabi--Yau pairs}
\label{sec:reduction}

We work over $\C$ until the final descent argument.  The reductions in
this section follow the strategy of Lazi\'c--Peternell
\cite[Sections 4 and 5]{LP}.  We use the terminating programs in
Theorem~\ref{input:models}, together with the abundance input
Theorem~\ref{input:abundance}, to make the precise reductions needed below.
The essential point is to retain control of two adjoints: an ordinary
semiample adjoint and a nearby generalized nef adjoint.

\subsection{The induction statements}

For a projective klt $\Q$-pair $(X,B)$ and a nef $\Q$-Cartier divisor $M$
on $X$, write
\[
                         T=K_X+B.
\]
Suppose that $T$ is pseudo-effective.  We will use the condition that,
for some $a\in\Q_{>0}$,
\begin{equation}\label{eq:full-condition}
 (T+aM)\cdot C>0
 \quad\text{for every integral curve $C$ through a very general point of $X$.}
\end{equation}
The divisor $T+aM$ need not be nef.  By
Lemma~\ref{lemma:pseudo-vg}, Condition~\eqref{eq:full-condition} holds
whenever $M$ has full nef dimension.  When $T+aM$ is nef, the condition
says exactly that its nef dimension is $\dim X$.

We prove the following two assertions together, by induction on $n$.
The notation $P_\sigma$ denotes the positive part of the divisorial
Zariski decomposition; its properties used in the proof are recalled
in Section~\ref{sec:positive}.
\begin{description}
\item[$(P_n)$] If $(X,B)$ is a projective klt $\Q$-pair of dimension $n$,
$T=K_X+B$ is pseudo-effective, $M$ is a nef $\Q$-Cartier divisor on $X$,
and \eqref{eq:full-condition} holds, then $T+cM$ is big for every
$c\in\Q_{>0}$.
\item[$(Z_n)$] If $(X,B)$ and $M$ have the same hypotheses, without
\eqref{eq:full-condition}, then there is a projective birational morphism
$w\colon W\to X$ from a smooth projective variety such that
\[
                         P_\sigma\bigl(w^*(T+M)\bigr)
\]
is a rational divisor numerically equivalent to a semiample
$\Q$-divisor on $W$.
\end{description}

Both assertions hold in dimensions zero and one.  In dimension one,
pseudo-effectivity of $T$ and nefness of $M$ mean that their degrees are
nonnegative.  Condition~\eqref{eq:full-condition} makes every positive
combination have positive degree, hence be ample.  A divisor of degree
zero is numerically equivalent to zero, and a divisor of positive degree
is ample.  These observations also give $(Z_1)$, since a nef divisor has
zero negative part.  On a point the assertions use the usual conventions.

Fix $n\geq2$, and assume $(P_j)$ and $(Z_j)$ for $j<n$.
At any stage we may replace $X$ by a small projective
$\Q$-factorialization.  The ordinary adjoint pulls back crepantly, $M$
pulls back to a nef divisor, and \eqref{eq:full-condition} persists by
using points in the isomorphism locus.  Bigness descends under this
birational pullback.  A smooth model proving $(Z_n)$ for the
$\Q$-factorialization also proves it for $X$.  We therefore impose
$\Q$-factoriality when running the programs below.

\subsection{Keeping a semiample adjoint fixed}

We first record the elementary ray comparison used in the second
program.  The nef part of a generalized pair in this argument is the
fixed nef $\Q$-Cartier b-divisor determined by $M$ on the original
variety.  Its trace on a later model need not be nef.

\begin{lemma}\label{lemma:generalized-ray-bound}
Let $(V,\Delta+\mathbf N)$ be a projective generalized klt $\Q$-pair
of dimension $n$, with effective boundary, $\Q$-factorial $V$, and nef
rational Cartier b-divisor data.  Every
$(K_V+\Delta+N_V)$-negative extremal ray has an integral rational curve
generator $C$ for which
\[
                -4n\leq (K_V+\Delta+N_V)\cdot C<0.
\]
The contraction, and the flip when it is small, can be realized using
an ordinary klt adjoint negative on the same ray.  Birational MMP
outputs remain $\Q$-factorial.  On a common resolution of a birational
step, the pullback of the input generalized adjoint equals the pullback
of its output transform plus an effective divisor exceptional over
the output.
\end{lemma}

\begin{proof}
Set $G=K_V+\Delta+N_V$, and fix a $G$-negative extremal ray.
For an ample rational divisor $A$ on $V$, choose $\eta>0$ rational and
small enough that $G+\eta A$ is still negative on the ray and
\[
                 -(G+\eta A)\cdot\gamma
                      \geq -\tfrac12G\cdot\gamma
\]
for its nonzero classes $\gamma$.  We may represent $G+\eta A$ as an
ordinary klt adjoint.  Here is the reason this remains valid when the
nef b-divisor does not descend to $V$.

Choose a projective log resolution $h\colon V'\to V$ carrying its nef
part $N'$ and write
\[
                  K_{V'}+\Delta'+N'=h^*G.
\]
The coefficients of $\Delta'$ are strictly less than one.  The
resolution can be chosen with an effective exceptional divisor $F$
such that $-F$ is $h$-ample.  For sufficiently small rational
$\varepsilon>0$, the divisor
$N'+\eta h^*A-\varepsilon F$ is ample.  A general effective rational
representative of this ample class can be chosen with sufficiently
small coefficients that, after adding $\varepsilon F$, it preserves
the sub-klt condition of $\Delta'$.  Its pushforward, added to
$\Delta$, is effective and defines a klt boundary $\Theta$ on $V$ with
\[
                       K_V+\Theta\sim_\Q G+\eta A.
\]
The crepant equality can be checked upstairs: its difference is
exceptional and relatively numerically trivial, hence zero by
negativity.

The ordinary cone theorem now gives a rational curve generator $C$
with $0<-(G+\eta A)\cdot C\leq2n$.  The preceding comparison yields
$-G\cdot C\leq4n$.  The contraction $c\colon V\to S$ is the
contraction of the same ray for this ordinary klt pair.  Since its
relative Picard number is one, there is a rational $\lambda>0$ such
that $G-\lambda(K_V+\Theta)$ is numerically trivial over $S$.
The Cartier descent statement for an ordinary klt extremal
contraction, after clearing denominators, gives a rational Cartier
divisor $D_S$ with
\[
                 G-\lambda(K_V+\Theta)\sim_\Q c^*D_S.
\]
In the small case, let $c^+\colon V^+\to S$ be the ordinary flip.
Taking strict transforms gives
\[
                 G^+\sim_\Q
                 \lambda(K_{V^+}+\Theta^+)+(c^+)^*D_S,
\]
which is ample over $S$.  Thus the ordinary flip is also the flip of
$G$.  The ordinary contraction and flip preserve $\Q$-factoriality
on the birational output.  Their effective exceptional adjoint
comparison, multiplied by $\lambda$, gives the asserted comparison
for $G$, since the terms pulled back from $S$ cancel.  With the nef
b-divisor fixed, this comparison is also the improvement of
generalized discrepancies used below.
\end{proof}

\begin{lemma}\label{lemma:two-nef-model}
Let $(X,B)$ be a projective $\Q$-factorial klt $\Q$-pair of dimension
$n$, let $T=K_X+B$ be pseudo-effective, and let $M$ be a nef
$\Q$-Cartier divisor on $X$.  Suppose
\eqref{eq:full-condition} holds.  Then there are a birational
contraction $X\dashrightarrow Y$, a number $s_0\in\Q_{>0}$, and a
common smooth projective resolution
\[
                       X\xleftarrow{p}W\xrightarrow{q}Y
\]
with the following properties.  If $B_Y$ and $N$ denote the transforms
of $B$ and $M$, and $U=K_Y+B_Y$, then
\begin{enumerate}
\item $(Y,B_Y)$ is klt and $\Q$-factorial, $U$ is semiample, and
$\kappa(Y,U)=\kappa(X,T)\geq0$;
\item $N$ is pseudo-effective, and $U+s_0N$ is nef of full nef
dimension;
\item there are effective $q$-exceptional rational divisors $E_0,E_1$
on $W$ such that
\begin{equation}\label{eq:two-comparisons}
 p^*T=q^*U+E_0,
 \qquad
 p^*(T+s_0M)=q^*(U+s_0N)+E_1.
\end{equation}
\end{enumerate}
\end{lemma}

\begin{proof}
By Theorem~\ref{input:models}, there is a terminating $T$-MMP, ending
on a $\Q$-factorial klt pair $(Y_0,B_0)$ with nef adjoint
$U_0=K_{Y_0}+B_0$.  Theorem~\ref{input:abundance} makes $U_0$
semiample.  The effective exceptional comparisons for the program
give
\[
                         \kappa(X,T)=\kappa(Y_0,U_0)\geq0.
\]
Let $N_0$ be the transform of $M$.  It is the trace on $Y_0$ of the
fixed nef b-divisor determined by $M$ on $X$.

Because this first program has finitely many steps, there is a
rational $s>0$ such that every one of them remains negative for the
transform of $T+uM$, for every rational $u\in(0,s]$.  Choose $s$ also
small enough to retain the strict sign on the flipped side of each
flip.  These are therefore programs for the corresponding generalized
klt pairs: they start generalized klt on $X$, and their discrepancies
improve at each step.  In particular, $(Y_0,B_0+s\mathbf M)$ is
generalized klt.  Pushforward to a $\Q$-factorial model preserves
pseudo-effectivity, so $N_0$ is pseudo-effective.

For completeness, this last numerical assertion applies also to
birational contractions with flips.  On a common resolution, numerical
pushforward to a $\Q$-factorial target is well defined: if a divisor
upstairs is numerically trivial, its difference from the pullback of
its pushforward is exceptional and relatively numerically trivial,
and the negativity lemma makes that difference zero.  Effective
approximation then proves preservation of pseudo-effectivity.

Choose $d>0$ so that $dU_0$ is Cartier and globally generated.  Choose
a rational $s_0\in(0,s)$ such that
\begin{equation}\label{eq:small-nef-coefficient}
                 \frac{1-s_0/s}{d}>4n\frac{s_0}{s}.
\end{equation}
The generalized adjoint $U_0+s_0N_0$ is pseudo-effective.  Apply
Theorem~\ref{input:models} to obtain a terminating program for it.
We claim that every step is $U_0$-trivial and is also negative for the
generalized adjoint with coefficient $s$, where throughout this
claim the symbols denote their transforms.

Suppose the claim holds up to the next step, and denote the current
transforms by $U_i,N_i$.  Then $dU_i$ is still Cartier and free, and
the pair with nef b-divisor coefficient $s$ is generalized klt.
For a $(U_i+s_0N_i)$-negative ray, the equality
\[
 U_i+s_0N_i=(1-s_0/s)U_i+(s_0/s)(U_i+sN_i)
\]
and nefness of $U_i$ show that the ray is $(U_i+sN_i)$-negative.
Choose its rational curve generator $C$ by
Lemma~\ref{lemma:generalized-ray-bound}.  If $U_i\cdot C>0$, then
integrality of $dU_i$ gives $U_i\cdot C\geq1/d$, whence
\[
 (U_i+s_0N_i)\cdot C
 \geq\frac{1-s_0/s}{d}-4n\frac{s_0}{s}>0,
\]
a contradiction.  Thus $U_i$ is trivial on the ray.

The morphism defined by $dU_i$ is constant on every fibre of the
extremal contraction and factors through it.  Its descended line
bundle is globally generated; pulling it to the flipped model, when
there is a flip, gives the transform of $dU_i$.  Hence that transform
is again Cartier and free, with equal pullbacks on a common
resolution.  Since the step is also negative for coefficient $s$,
generalized klt singularities with that coefficient persist.  This
proves the claim by induction through the finite program.

Let $Y$ be its output.  The ordinary adjoint $U$ is crepant throughout
the second program, so $(Y,B_Y)$ is klt, $U$ is semiample, and its
Iitaka dimension equals that of $T$.  The generalized adjoint
$U+s_0N$ is nef.  Combining the pullback comparisons of the two
programs gives \eqref{eq:two-comparisons} with effective
$q$-exceptional divisors.  The transform $N$ remains pseudo-effective.

It remains to prove full nef dimension.  Suppose instead that the
nef reduction of $U+s_0N$ has a positive-dimensional very general
compact fibre $F$.  Shrink the base of the almost holomorphic map so
that these compact fibres lie entirely in its morphism locus.  The
image in $Y$ of the exceptional locus of $q$ has
codimension at least two.  A general fibre $F$ meets each of its
components in codimension at least two, or not at all: discard the
images of components that do not dominate the nef-reduction base,
and use the fibre dimension theorem for the others.  We may
therefore choose a curve $C\subset F$ through a very general point,
avoiding this exceptional image.  If $\dim F>1$, take sufficiently
general complete intersections in $F$ through that point; if
$\dim F=1$, take $F$ itself.

Choose the point also in the common isomorphism locus of the
birational contraction and with the very general conditions of
\eqref{eq:full-condition} transported from $X$.  Since $U$ is nef
and $N$ is pseudo-effective, Lemma~\ref{lemma:pseudo-vg} gives
\[
                    U\cdot C=N\cdot C=0,
\]
because their positive combination has degree zero on $F$.
The strict lift $\widetilde C\subset W$ avoids $E_0$ and $E_1$.
Both equalities in \eqref{eq:two-comparisons} then give
\[
               p^*T\cdot\widetilde C
                =p^*M\cdot\widetilde C=0.
\]
The curve $p(\widetilde C)$ passes through the chosen very general
point of $X$ and contradicts \eqref{eq:full-condition}.
\end{proof}

\subsection{Positive Iitaka dimension and the remaining case}

The preceding construction leaves only one obstruction to bigness:
the ordinary semiample adjoint may define a map to a point.  If it
defines a positive-dimensional base, the induction hypothesis handles
its fibres.

\begin{proposition}\label{prop:full-positive}
Assume $(P_j)$ for $j<n$.  Let $(X,B)$ be a projective klt
$\Q$-pair of dimension $n$, let $T=K_X+B$ be pseudo-effective, and
let $M$ be a nef $\Q$-Cartier divisor on $X$.  If
\eqref{eq:full-condition} holds and $\kappa(X,T)>0$, then $T+cM$ is
big for every $c\in\Q_{>0}$.
\end{proposition}

\begin{proof}
After a small $\Q$-factorialization, apply
Lemma~\ref{lemma:two-nef-model}.  Let $h\colon Y\to S$ be the
contraction defined by the semiample divisor $U$.  There is an ample
rational divisor $A_S$ on $S$ with $U\sim_\Q h^*A_S$, and
$\dim S=\kappa(X,T)>0$.  Put $Q=U+s_0N$.

If a very general fibre $F$ has positive dimension, it has dimension
strictly less than $n$.  The restriction pair $(F,B_Y|_F)$ is klt
and
\[
                    K_F+B_Y|_F\sim_\Q U|_F\sim_\Q0.
\]
Moreover, $Q|_F$ is nef of full nef dimension.  Indeed, a very general
point of a very general $F$ can be chosen outside the countable union
excluded in the full-nef-dimension assertion for $Q$ on $Y$.
Since $U|_F\equiv0$, the divisor $N|_F$ is nef.  Apply $(P_{\dim F})$
to this restriction pair and $s_0N|_F$ to conclude that $Q|_F$ is big.
When the fibre dimension is zero, relative bigness is automatic.

Thus $Q$ is big over $S$.  To justify the passage to the generic
fibre, fix an ample Cartier divisor $H$ on $Y$ and shrink $S$ so
that $h$ is flat.  For every positive integer $\ell$ clearing the
denominators of $Q$, upper semicontinuity makes the locus
\[
        \bigl\{s\in S:
        H^0(F_s,\OO_{F_s}((\ell Q-H)|_{F_s}))\ne0\bigr\}
\]
closed.  Bigness on a very general fibre puts its base point in one
of these countably many loci.  Choose that point outside every
locus that is proper; the locus containing it must then be all of
$S$.  On the generic fibre, $\ell Q$ is consequently the sum of
the ample divisor $H$ and an effective divisor, which proves the
asserted relative bigness.  A relatively big divisor becomes big after adding
a sufficiently large pullback of an ample divisor on the base;
equivalently, apply the usual relative form of Kodaira's lemma
\cite[Lemma 3.2.1]{BCHM}.  It follows that
\[
                         Q+bU=(1+b)U+s_0N
\]
is big for some rational $b>0$.  Equation~\eqref{eq:two-comparisons}
then shows that $(1+b)T+s_0M$ is big.

Finally, fix $c>0$ rational and choose $\lambda>0$ rational with
$\lambda(1+b)<1$ and $\lambda s_0<c$.  Then
\[
 T+cM=\lambda\bigl((1+b)T+s_0M\bigr)
       +\bigl(1-\lambda(1+b)\bigr)T+(c-\lambda s_0)M.
\]
The first summand is big and the other two are pseudo-effective.
Their sum is big, as required.
\end{proof}

\begin{proposition}\label{prop:reduction-cy}
Assume $(P_j)$ for $j<n$.  Suppose that on every projective klt
$\Q$-pair $(V,\Delta)$ of dimension $n$ with
$K_V+\Delta\equiv0$, every nef $\Q$-Cartier divisor of full nef
dimension is big.  Then $(P_n)$ holds.
\end{proposition}

\begin{proof}
Use a small $\Q$-factorialization and
Lemma~\ref{lemma:two-nef-model}.  The nonnegative Iitaka dimension
of $T$ is either positive, in which case
Proposition~\ref{prop:full-positive} applies, or zero.  In the latter
case, semiampleness of $U$ implies $U\sim_\Q0$.  The pair $(Y,B_Y)$
is therefore a klt Calabi--Yau pair.  The nef divisor $U+s_0N$ has
full nef dimension, so it is big by the stated hypothesis.
The second equality in \eqref{eq:two-comparisons} gives bigness of
$T+s_0M$.  The positive-combination argument at the end of
Proposition~\ref{prop:full-positive} gives bigness of $T+cM$ for
every rational $c>0$.
\end{proof}

Consequently, the remaining work for $(P_n)$ is the full-nef-dimension
statement on klt Calabi--Yau pairs.  Theorem~\ref{thm:terminal-big}
establishes its terminal, zero-boundary form.  We then pass from that
form to all klt Calabi--Yau pairs in Section~\ref{sec:completion}.
